\section{Csoport-orbitális konzisztencia a zenei VAE-modellekben}

\subsection{Motiváció és elméleti háttér}

A zenei struktúrák egyik legfontosabb sajátossága az ismétlődő mintázatok és szimmetrikus szerkezetek jelenléte \cite{Temperley2001, Lerdahl2004}.  
Egy adott motívum felismerhető marad akkor is, ha transzponálva, inverzióval, vagy más zenei művelettel módosítjuk — ez a jelenség a \textit{hangnemi invariancia}.  
A legtöbb mélytanulási modell, különösen a klasszikus Variációs Autoencoder (VAE) \cite{Kingma2013}, nem kezeli explicit módon ezt az invarianciát: minden variánst külön mintaként tanul, ezzel redundanciát és információveszteséget okozva.

Ebből adódott a kutatás vezérkérdése:  
\begin{quote}
\textit{Hogyan építhető be a zene szimmetriacsoportjainak szerkezete egy VAE-modell tanulási folyamatába oly módon, hogy a latens tér invariáns legyen a zenei transzformációkkal szemben?}
\end{quote}

\subsection{Alapötlet}

Az alapötlet az, hogy minden zenei motívumhoz nemcsak önmagát, hanem az általa generált teljes \emph{csoportorbitot} (\(G \cdot x\)) is figyelembe vesszük.  
Ahelyett, hogy a modell minden orbit-elemhez külön latens vektort rendelne, az encoder-t úgy tanítjuk, hogy azok azonos reprezentációt eredményezzenek.  
Ezt a jelenséget \textbf{csoport-orbitális konzisztenciának} (Group-Orbit Latent Consistency, GOLC) nevezzük.

Formálisan: ha \(G\) egy transzformációs csoport (pl. hangnemi transzpozíciók), akkor az encodernek teljesülnie kell a következő feltételnek:
\[
Enc_\phi(g \cdot x) \approx Enc_\phi(x), \quad \forall g \in G.
\]
Ez biztosítja, hogy a latens térben a zenei struktúrák \emph{formája} megmarad, míg a felszíni változások (pl. hangmagasság-eltolások) eltűnnek.  
Ennek következményeként a modell \emph{strukturálisan generalizál} — nem konkrét dallamokat, hanem zenei viszonyokat tanul.

\subsection{A gondolatmenet kialakulása}

A VAE-architektúra természetes módon képes latens absztrakciókat tanulni, azonban ezek nincsenek kötve semmilyen szimmetriafeltételhez.  
A zene esetében ez különösen hátrányos, hiszen a kompozíciós logika nagyrészt a transzformációkon alapul (transzpozíció, augmentáció, retrográd, stb.).  
A csoportelmélet \cite{Armstrong1988} kínálta a legtermészetesebb formalizmust ennek az invarianciának a modellezésére.

A fejlesztés során felismerésre került, hogy a VAE encoder és decoder párosa egy nemlineáris leképezésrendszerként felfogható, ahol az encoder egy \(\mathcal{X}\) térből egy \(\mathcal{Z}\) latens térbe képez.  
Ha a \(\mathcal{X}\)-en egy csoport \(G\) hat, akkor a cél az, hogy az \(\mathcal{Z}\)-ben ez a hatás \emph{invariánssá} váljon:
\[
Enc_\phi(g \cdot x) = Enc_\phi(x), \quad \forall g \in G.
\]
Ezt az invarianciát regularizáló taggal vezetjük be a veszteségfüggvénybe.

\subsection{Formális definíciók}

\begin{definition}[Orbitális konzisztencia]
Egy encoder \(Enc_\phi\) orbitálisan konzisztens, ha bármely \(x \in \mathcal{X}\) és \(g \in G\) esetén:
\[
Enc_\phi(g \cdot x) = Enc_\phi(x).
\]
\end{definition}

\begin{definition}[Kanonikus reprezentáció]
Az orbit reprezentációinak átlaga:
\[
z_c(x) = \frac{1}{|G|} \sum_{g\in G} Enc_\phi(g\cdot x).
\]
\end{definition}

\begin{definition}[Orbitális konzisztencia veszteség]
\[
L_{\text{orbit}} = \frac{1}{|G|} \sum_{g\in G} \|Enc_\phi(g\cdot x) - z_c(x)\|^2.
\]
\end{definition}

\begin{definition}[Teljes veszteségfüggvény]
\[
L_{\text{total}} = L_{\text{recon}} + \beta_{\text{KL}} L_{\text{KL}} + \beta_{\text{orbit}} L_{\text{orbit}}.
\]
\end{definition}

\subsection{Formális igazolások}

\begin{proposition}[Orbitális invariancia feltétele]
Ha \(L_{\text{orbit}} = 0\), akkor az encoder invariáns az orbiton.
\end{proposition}
\begin{proof}
A \(L_{\text{orbit}}\) tag definíció szerint nemnegatív.  
A nullpontban minden orbit-elem reprezentációja azonos a kanonikus vektorral, tehát az encoder invariáns a csoporthatásra.
\end{proof}

\begin{proposition}[Posterior disztribúciós stabilitás]
Ha az encoder posteriorjai normáleloszlásúak \(q_\phi(z|x) = \mathcal{N}(\mu_x, \Sigma_x)\), és \(L_{\text{orbit}}\) kiterjed mindkét paraméterre, akkor:
\[
\mu_{g\cdot x} = \mu_x, \quad \Sigma_{g\cdot x} = \Sigma_x, \quad \forall g\in G.
\]
\end{proposition}
\begin{proof}
A Frobenius-norm minimalizálása konvex a kovarianciákra, így a minimumot az egyenlőség adja.  
Ez ekvivalens azzal, hogy a posterior eloszlások invariánsak a csoporthatásra \cite{Kondor2018, Cohen2016}.
\end{proof}

\begin{proposition}[Tanulási stabilitás]
Az orbitális regularizáció csökkenti a \(L_{\text{KL}}\) gradiensének orbiton belüli varianciáját:
\[
Var_G(\nabla_\phi L_{\text{KL}}) \to 0.
\]
\end{proposition}
\begin{proof}
Ha minden orbit-elem ugyanazt a latens eloszlást adja, a veszteségfüggvény gradiensének értéke minden transzformációra azonos lesz.  
Ez kiegyenlíti a gradiens-fluktuációkat, gyorsabb és stabilabb konvergenciát eredményezve, hasonlóan a data-augmentation regularizáció hatásához \cite{Hadjeres2017}.
\end{proof}

\begin{proposition}[Dekódolási kompatibilitás]
Ha \(Enc_\phi\) invariáns és \(Dec_\theta\) folytonos, akkor
\[
Dec_\theta(Enc_\phi(g\cdot x)) \approx g\cdot Dec_\theta(Enc_\phi(x)).
\]
\end{proposition}
\begin{proof}
Az encoder invarianciája miatt a dekóder bemenete invariáns a csoporthatásra.  
A folytonosság biztosítja, hogy a dekódolási térben a kimenetek közel equivariánsak legyenek a csoporttranszformációkra \cite{Cohen2016}.
\end{proof}

\subsection{Empirikus következmények és várakozások}

Elméleti megfontolásaink alapján a GOLC bevezetése:

\begin{itemize}
    \item \textbf{Csökkenti a redundanciát}, hiszen azonos orbit-elemeket egyetlen reprezentációba tömörít.
    \item \textbf{Javítja az általánosítást}, mert a modell zenei transzformációkra invariáns struktúrát tanul.
    \item \textbf{Stabilizálja a tanulást}, mivel a regularizáció kisimítja a latens tér topológiáját.
    \item \textbf{Zeneileg interpretálhatóbb} reprezentációt hoz létre, mivel a latens tér dimenziói közvetlenül megfeleltethetők zenei transzformációknak.
\end{itemize}

\subsection{Összehasonlítás a hagyományos VAE-modellel}

% IDE SZÚRD BE: az elemzést, ahol a régi VAE modelledet és a GOLC-VAE változatot összeveted.

\subsection{Metrikák és javaslatok az összehasonlításhoz}

% IDE: Mérések és összehasonlítási javaslatok
% Ajánlott metrikák:
% 1. Orbit Latent Distance (OLD): az orbiton belüli reprezentációk átlagos euklideszi távolsága.
% 2. KL Variance: a KL-divergencia orbiton belüli szórása.
% 3. Reconstruction MSE: a rekonstrukciós hiba.
% 4. Pitch-Class Entropy: hangnem-stabilitás mérése.
% 5. Tonality Preservation Score: mennyire marad konzisztens a zenei szerkezet transzformáció után.
% Ezekhez grafikonokat vagy táblázatokat érdemes beilleszteni.

\subsection{Összegzés}

A csoport-orbitális konzisztencia elve a VAE-architektúrák természetes kiterjesztése, amely a zene transzformációs szimmetriáit explicit módon integrálja a tanulási folyamatba.  
A módszer elméletileg igazolt, gyakorlatban pedig várhatóan javítja a zenei kohéziót, csökkenti az adatigényt és növeli a modell robusztusságát.

\subsection{Hivatkozások}

\bibliographystyle{plain}
\bibliography{references}
